\section{Pruebas unitarias}
\label{sec:pruebas}

El enunciado exige documentar al menos dos pruebas unitarias por método.
La sección ``Pruebas unitarias'' de
\texttt{notebooks/tp1\_optimizacion.ipynb} implementa exactamente dos
pruebas para cada algoritmo. Cada prueba es una función documentada con
Doxygen que falla mediante \texttt{assert} o
\texttt{torch.testing.assert\_close} si el resultado no coincide con el
esperado. En la ejecución reportada las seis pruebas se superaron (el
notebook imprime ``Seis pruebas unitarias superadas: dos de GD, dos de
RMSProp y dos de PSO'').

\subsection{Descenso del gradiente}

\begin{description}
    \item[\texttt{probar\_actualizacion\_gd}] Ejecuta un paso de GD desde
    $(1,-2)$ sobre $f_0$ con $\alpha=0.1$. Como $\nabla f_0(x)=2x$, el punto
    esperado es $x_1=x_0-\alpha\cdot2x_0=(0.8,-1.6)$. La prueba compara el
    punto inicial y el actualizado del historial con los esperados, y los
    valores de función con $f_0(x_0)$ y $f_0(x_1)$. Verifica que el
    gradiente de Autograd y la regla de actualización coinciden con la
    derivada analítica.
    \item[\texttt{probar\_historiales\_gd}] Ejecuta 5 pasos desde $(4,-3)$
    sobre $f_0$ con $\alpha=0.1$. La solución analítica es
    $x_k=0.8^k\,x_0$. La prueba comprueba que los historiales tienen forma
    $(6,2)$ y $(6,)$ (punto inicial más 5 pasos) y que todos los puntos y
    valores coinciden con la solución cerrada. Verifica que el historial
    completo, y no solo el último paso, es correcto.
\end{description}

\subsection{RMSProp}

\begin{description}
    \item[\texttt{probar\_actualizacion\_rmsprop}] Ejecuta un único paso de
    RMSProp desde $(1,-2)$ sobre $f_0$, con $\alpha=0.1$, $\gamma=0.9$ y
    $\epsilon=10^{-6}$, y compara el punto y el valor de función
    resultantes contra los calculados explícitamente con
    $s_1=(1-\gamma)g_0^2$ y $x_1=x_0-\alpha g_0/\sqrt{s_1+\epsilon}$, usando
    \texttt{torch.testing.assert\_close}. Verifica que la implementación
    vectorizada coincide con la fórmula matemática.
    \item[\texttt{probar\_historiales\_rmsprop}] Ejecuta 25 iteraciones de
    RMSProp sobre $f_0$, $f_1$ y $f_2$ desde $(4,-3)$, $(2.5,-1.5)$ y
    $(6,6)$, con $\alpha=0.5$, $0.05$ y $0.4$ y $\gamma=0.9$, y comprueba
    para cada caso que: el historial de puntos tiene forma $(26,2)$ y el de
    valores $(26,)$; todos los valores son finitos; el primer punto coincide
    con el punto inicial; y el valor final es estrictamente menor que el
    inicial.
\end{description}

\subsection{Enjambre de partículas}

\begin{description}
    \item[\texttt{probar\_actualizacion\_pso}] Ejecuta 4 iteraciones con 5
    partículas sobre $f_0$ con $c_1=c_2=0$ (semilla 7). Sin atracción, la
    ecuación de velocidad se reduce a $v\leftarrow v$, así que cada
    partícula debe moverse en línea recta con velocidad constante:
    $x^{(k)}=x^{(0)}+k\,(x^{(1)}-x^{(0)})$. La prueba verifica esta relación
    para todas las partículas e iteraciones, y que el mejor global inicial
    sea la partícula de menor $f_0$ en $t=0$, con su valor.
    \item[\texttt{probar\_historiales\_pso}] Ejecuta 25 iteraciones con 20
    partículas, $c_1=c_2=0.5$ y semillas fijas sobre $f_0$, $f_1$ y $f_2$.
    Comprueba las formas de los historiales ($(26,2)$, $(26,)$ y
    $(26,20,2)$), que todos los valores sean finitos, que la primera
    partícula empiece en el punto compartido, que el valor del mejor global
    en cada iteración sea exactamente el mínimo acumulado
    (\texttt{torch.cummin}) de $f$ sobre todas las posiciones visitadas por
    el enjambre, y que $f$ evaluada en las posiciones del mejor global
    coincida con los valores reportados.
\end{description}
